\begin{abstract}
Hospital-associated saliva deposition can seed high-risk pathogens across shared surfaces. This work studies unsupervised segmentation of saliva deposits illuminated under UV-C, where only ten color images are available and ground truth masks only cover four centralized patches per frame. We explore a hierarchy of methods: a MATLAB-based luminance/thresholding pipeline (Method A), an exploratory k-means segmentation (Method C), and our probabilistic Gaussian mixture model (Method D) that reuses k-means centers, runs ten EM initializations per image, and selects the most likely fit by pseudo-log-likelihood. The resulting GMM probability maps are binarized with morphological cleanup and evaluated on accuracy, precision, recall, and Dice; the \(K{=}4\) model achieves the highest Dice while maintaining recall above 0.96. Our comparison highlights how converting hard clusters into likelihood-ranked mixtures improves recall while still recovering reasonable precision, providing a practical unsupervised pipeline for low-data airborne contamination imagery.
\end{abstract}